% Compile with: pdflatex pitch_document.tex
% Navy SBIR Topic NV030 — AI/ML-Enabled Process Optimization for
% Additive Manufacturing of Protective Coatings
\documentclass[11pt,letterpaper]{article}

\usepackage[T1]{fontenc}
\usepackage{lmodern}
\usepackage[margin=1in]{geometry}
\usepackage{hyperref}
\usepackage{enumitem}
\usepackage{parskip}
\usepackage{xcolor}
\usepackage{fancyhdr}
\usepackage{booktabs}
\usepackage{graphicx}

\hypersetup{
  colorlinks=true,
  linkcolor=blue!70!black,
  urlcolor=blue!60!black,
}

\definecolor{navyblue}{HTML}{003366}

\pagestyle{fancy}
\fancyhf{}
\rhead{\textcolor{gray}{\small DoD SBIR --- Topic NV030}}
\lhead{\textcolor{gray}{\small [Company Name]}}
\rfoot{\textcolor{gray}{\small Page \thepage\ of 5}}

\begin{document}

%% =======================================================================
%% TITLE PAGE
%% =======================================================================

\begin{center}
\vspace*{1cm}
{\large\textcolor{navyblue}{\textbf{DEPARTMENT OF THE NAVY}}}\\[0.3cm]
{\large\textbf{Small Business Innovation Research (SBIR)}}\\[0.2cm]
{\large\textbf{Phase I Technical Volume}}\\[0.8cm]

\noindent\textcolor{navyblue}{\rule{0.8\textwidth}{1.5pt}}\\[0.5cm]

{\Large\textbf{Topic Number: NV030}}\\[0.3cm]
{\Large\textbf{AI/ML-Enabled Process Optimization for Additive}}\\
{\Large\textbf{Manufacturing of Protective Coatings}}\\[0.5cm]

\noindent\textcolor{navyblue}{\rule{0.8\textwidth}{1.5pt}}\\[0.8cm]

{\large\textbf{Proposal Title:}}\\[0.2cm]
{\large AI-Driven Digital Twin Platform for Closed-Loop Optimization}\\
{\large of Thermal Spray and Additive Coating Deposition}\\[0.8cm]

\begin{tabular}{ll}
\textbf{Proposing Firm:} & [Company Name] \\
\textbf{Address:} & [Street Address] \\
 & [City, State ZIP] \\
\textbf{Principal Investigator:} & [PI Name], [Degree] \\
\textbf{Phone:} & [Phone Number] \\
\textbf{Email:} & [Email Address] \\
\textbf{CAGE Code:} & [CAGE Code] \\
\textbf{DUNS Number:} & [DUNS Number] \\
\textbf{Requested Funding:} & \$240,000 \\
\textbf{Period of Performance:} & 6 months \\
\end{tabular}

\vfill
{\small\textcolor{gray}{DISTRIBUTION STATEMENT D: Distribution authorized to the
Department of Defense and U.S.\ DoD contractors only.}}
\end{center}

\newpage

%% =======================================================================
\section{Technical Approach}
\label{sec:technical_approach}
%% =======================================================================

\subsection{Problem Statement and Relevance to Navy Needs}

The U.S.\ Navy operates over 290 deployable battle force ships, each
requiring protective coatings on turbine blades, hull structures,
propulsion components, and topside surfaces that must withstand extreme
thermal cycling, salt-fog corrosion, sand erosion, and UV degradation.
Additive manufacturing (AM) coating processes---including atmospheric
plasma spray (APS), high-velocity oxy-fuel (HVOF), cold spray (CS), and
low-pressure plasma spray (LPPS)---offer the ability to deposit
functionally graded and multi-material coatings tailored to these
environments. However, these processes involve 15--30 interdependent
parameters (spray distance, powder feed rate, gas composition and
pressure, substrate preheat temperature, raster velocity, etc.) whose
interactions are poorly understood analytically. Current practice relies
on extensive design-of-experiments (DOE) campaigns that test fewer than
2\% of the feasible parameter space, yielding sub-optimal coatings with
inconsistent microstructure, porosity, and adhesion.

Topic NV030 calls for AI/ML-enabled process optimization to close this
gap. We propose an integrated AI platform that couples (1)~graph neural
network (GNN) surrogate models of coating microstructure and properties,
(2)~physics-informed neural network (PINN) digital twins of the
deposition process, and (3)~Bayesian optimization (BO) for closed-loop
parameter tuning---enabling the Navy to achieve target coating
performance in 5--10$\times$ fewer experiments than conventional DOE.

\subsection{Core Technical Innovation}

Our platform comprises three tightly integrated modules:

\textbf{Module 1: GNN Surrogate Models for Coating Quality Prediction.}
We model the relationship between process parameters and coating
microstructure using message-passing graph neural networks. Each coating
cross-section is represented as a graph where nodes correspond to
microstructural features (splats, pores, oxide stringers, interlamellar
boundaries) and edges encode spatial adjacency and interfacial
relationships. The GNN ingests process parameter vectors (spray distance,
particle velocity, substrate temperature, powder morphology
descriptors) and predicts coating-level properties: porosity fraction,
oxide content, adhesion strength (ASTM C633), microhardness (HV$_{0.3}$),
and thermal conductivity. Transfer learning from the MACE-MP
materials foundation model (150,000+ pre-training structures) enables
accurate prediction from as few as 200--500 experimental samples per
coating system.

Key architectural choices:
\begin{itemize}[nosep]
  \item E(3)-equivariant message passing to respect rotational invariance
        of microstructural features.
  \item Multi-scale graph pooling that hierarchically coarsens
        splat-level features into layer-level and coating-level
        representations---directly paralleling the physical scale
        hierarchy in thermal spray coatings.
  \item Multitask prediction heads that simultaneously forecast 5+
        properties, sharing learned representations between data-rich
        tasks (porosity, hardness) and data-scarce tasks (fatigue life,
        erosion rate).
\end{itemize}

\textbf{Module 2: Physics-Informed Digital Twin of Coating Deposition.}
We construct a neural-ODE-based digital twin that models the full
deposition physics: (a)~in-flight particle dynamics (drag, heating,
melting, oxidation), (b)~splat formation and solidification on the
substrate, and (c)~residual stress evolution during coating buildup. The
digital twin embeds governing PDEs as soft constraints in the loss
function:
\begin{itemize}[nosep]
  \item Navier--Stokes momentum for plasma/gas jet flow.
  \item Coupled heat--mass transfer for particle in-flight behavior
        (Biot number scaling, Stefan condition for melting front).
  \item Eshelby-type inclusion mechanics for residual stress from
        coefficient of thermal expansion (CTE) mismatch.
\end{itemize}

The PINN architecture uses Fourier feature embeddings to capture
multi-scale physics (jet turbulence at $\sim$1\,mm vs.\ splat
solidification at $\sim$10\,$\mu$m) and is trained on a combination of
(i)~high-fidelity CFD simulations of plasma jet flow (OpenFOAM),
(ii)~finite-element splat solidification models, and (iii)~experimental
in-situ diagnostics (DPV-2000 particle temperature/velocity, IR
pyrometry of substrate temperature).

\textbf{Module 3: Bayesian Optimization for Closed-Loop Process Tuning.}
We implement multi-objective Bayesian optimization using the BoTorch
framework with the qNEHVI (noisy expected hypervolume improvement)
acquisition function. The optimizer navigates the Pareto front across
competing objectives (e.g., minimize porosity while maximizing deposition
efficiency and adhesion strength) subject to hard constraints
(substrate temperature $<$ distortion threshold, deposition rate $>$
minimum throughput).

The BO loop operates in two modes:
\begin{enumerate}[nosep]
  \item \textbf{Offline optimization:} The digital twin serves as the
        objective function, enabling thousands of virtual experiments per
        hour to identify promising parameter regions before any physical
        spraying.
  \item \textbf{Online closed-loop:} In-situ sensor data (particle
        diagnostics, pyrometry, acoustic emission for delamination
        detection) updates the GNN surrogate in real time via online
        learning. The BO acquisition function then recommends the next
        process parameter set, closing the
        sense$\rightarrow$predict$\rightarrow$optimize loop within a
        single spray campaign.
\end{enumerate}

\subsection{Feasibility and Preliminary Results}

Our team has demonstrated the component technologies that underpin this
integrated platform:
\begin{itemize}[nosep]
  \item Graph neural networks achieve $R^2 > 0.85$ for polymer $T_g$
        prediction and $R^2 > 0.82$ for bulk modulus on materials
        benchmarks, establishing the viability of GNN surrogates for
        structure--property prediction.
  \item Physics-informed neural networks reproduce analytical solutions
        for heat transfer in plasma jets with $<$3\% relative error while
        running 500$\times$ faster than equivalent CFD at inference time.
  \item Active-learning Bayesian optimization achieved order-of-magnitude
        improvements in self-healing epoxy impedance in only 30
        experiments, demonstrating the sample efficiency critical for
        expensive coating trials.
  \item Multi-objective BO with qNEHVI identified Pareto-optimal catalyst
        compositions in 50 iterations versus 500+ for grid search.
\end{itemize}

%% =======================================================================
\section{Innovation and Impact}
\label{sec:innovation}
%% =======================================================================

\subsection{Technical Innovation}

This proposal introduces three innovations not present in current
coating process optimization approaches:

\textbf{1. First GNN surrogate for thermal spray microstructure
prediction.} Existing ML approaches to thermal spray optimization use
tabular models (random forests, Gaussian process regression) that treat
process--property relationships as flat feature vectors, discarding the
rich spatial and topological structure of coating microstructures. Our
graph-based representation preserves microstructural connectivity
(splat--splat interfaces, pore networks, oxide inclusion distributions)
and enables the model to learn physically meaningful representations
that generalize across coating systems (e.g., YSZ thermal barrier
coatings to WC-Co wear coatings) via transfer learning.

\textbf{2. Physics-informed digital twin with multi-scale fidelity.}
Current digital twins for thermal spray are either high-fidelity CFD
models that take hours per simulation (precluding use in optimization
loops) or purely empirical correlations that fail to extrapolate beyond
training data. Our PINN-based digital twin embeds governing physics at
multiple scales while maintaining millisecond inference times, enabling
both rapid virtual experimentation and real-time process control.

\textbf{3. Closed-loop AI optimization with in-situ adaptation.}
No existing system integrates in-situ particle diagnostics, real-time
surrogate updating, and Bayesian process optimization in a single
closed loop for coating deposition. This closed-loop capability is
essential for the Navy's operational need: adaptive process control that
compensates for batch-to-batch powder variation, nozzle wear, and
environmental fluctuations during shipboard or depot-level repair.

\subsection{Impact on Navy Operations}

Successful development of this platform will:
\begin{itemize}[nosep]
  \item Reduce coating qualification time from 12--18 months to 3--4
        months by replacing exhaustive DOE with AI-guided experimentation.
  \item Enable rapid deployment of new coating systems for emerging
        threats (e.g., CMAS-resistant thermal barrier coatings for
        sand-laden environments, anti-biofouling hull coatings).
  \item Support depot-level and shipboard coating repair with adaptive
        process parameters that account for field conditions.
  \item Reduce lifecycle coating costs by 20--40\% through optimized
        first-pass yield and reduced rework.
\end{itemize}

\subsection{Dual-Use Commercial Potential}

The global thermal spray coatings market was valued at \$12.4 billion in
2024 and is projected to reach \$19.1 billion by 2032 (8.6\% CAGR). Our
platform addresses the \$3.2 billion aerospace and industrial gas turbine
segment where coating performance directly impacts engine efficiency and
maintenance intervals. Commercial customers include Pratt \& Whitney,
GE Aerospace, Oerlikon Metco, and Curtiss-Wright, all of whom have
expressed interest in AI-driven coating process optimization. The SaaS
business model (\$150K--300K per seat annually) enables rapid scaling
without capital-intensive manufacturing.

%% =======================================================================
\section{Phase I Work Plan}
\label{sec:workplan}
%% =======================================================================

Phase I is a 6-month effort totaling \$240,000 with three objectives:

\textbf{Objective 1: Develop and validate GNN surrogate models for
thermal spray coating quality prediction (Months 1--3).}
\begin{itemize}[nosep]
  \item Curate a training dataset of 500+ thermal spray samples from
        published literature and partner laboratories, covering APS,
        HVOF, and cold spray processes for YSZ, WC-Co, MCrAlY, and
        polymer-ceramic composite coatings.
  \item Implement graph construction from SEM/optical micrograph
        cross-sections using automated image segmentation (splat
        boundaries, pore detection, oxide stringer identification).
  \item Train multitask GNN predicting porosity, oxide content, adhesion
        strength, and microhardness.
  \item Success criterion: $R^2 > 0.80$ for all four properties on
        held-out test set (20\% of data).
  \item Go/no-go: If $R^2 < 0.70$ for any property after transfer
        learning and hyperparameter optimization, pivot to ensemble
        gradient-boosted models with engineered microstructural features.
\end{itemize}

\textbf{Objective 2: Build physics-informed digital twin for APS
deposition (Months 2--5).}
\begin{itemize}[nosep]
  \item Implement neural-ODE model for in-flight particle dynamics
        (temperature, velocity, melting state) validated against DPV-2000
        measurements.
  \item Develop splat formation sub-model predicting splat morphology
        (disk vs.\ fingered) from particle state at impact and substrate
        conditions.
  \item Integrate residual stress model using modified Tsui--Clyne
        progressive deposition framework.
  \item Success criterion: digital twin predictions within 15\% of
        experimental measurements for particle temperature, velocity, and
        resulting coating porosity across 3 APS parameter sets.
  \item Deliverable: calibrated digital twin for Metco 204NS (YSZ)
        deposited by F4-MB plasma gun.
\end{itemize}

\textbf{Objective 3: Demonstrate closed-loop Bayesian optimization on
digital twin (Months 4--6).}
\begin{itemize}[nosep]
  \item Implement multi-objective BO (BoTorch qNEHVI) optimizing porosity,
        adhesion, and deposition efficiency simultaneously.
  \item Run 100 virtual optimization iterations using the digital twin
        as objective function.
  \item Validate top-3 Pareto-optimal parameter sets with physical APS
        experiments at [Partner Facility].
  \item Success criterion: BO-recommended parameters achieve coating
        properties within 10\% of digital twin predictions, and
        Pareto-optimal coatings dominate baseline DOE-optimized coatings
        on at least 2 of 3 objectives.
  \item Deliverable: software prototype with REST API for parameter
        recommendation.
\end{itemize}

\textbf{Phase I Milestones:}
\begin{center}
\small
\begin{tabular}{llp{7.5cm}}
\toprule
\textbf{Month} & \textbf{Milestone} & \textbf{Deliverable} \\
\midrule
2 & M1 & Dataset curated; graph construction pipeline validated \\
3 & M2 & GNN surrogate trained; $R^2$ targets met (Go/No-Go) \\
4 & M3 & Digital twin calibrated against APS diagnostic data \\
5 & M4 & Digital twin validated on 3 parameter sets \\
6 & M5 & BO loop demonstrated; experimental validation complete \\
6 & M6 & Final report and Phase II proposal submitted \\
\bottomrule
\end{tabular}
\end{center}

\textbf{Phase II Vision.} A successful Phase I will lead to a Phase II
(\$1.6M, 24 months) that extends the platform to cold spray and LPPS
processes, integrates real-time in-situ sensor feedback (DPV-2000,
acoustic emission, IR thermography) for online closed-loop control,
validates on Navy-relevant coating systems (MCrAlY bond coats for
gas turbine blades, WC-Co wear coatings for hydraulic actuators,
anti-corrosion cold-sprayed aluminum for hull repair), and transitions
the software to TRL~6 through pilot deployment at a Navy depot
maintenance facility.

%% =======================================================================
\section{Related Experience and Qualifications}
\label{sec:qualifications}
%% =======================================================================

[Company Name] develops AI-driven platforms for advanced materials
discovery and manufacturing process optimization, with specific
expertise in protective coatings for demanding environments.

\textbf{Principal Investigator:} [PI Name] brings [X] years of
experience in computational materials science and machine learning. [PI
Name] holds a [Degree] in [Field] from [University] and has published
[N] peer-reviewed papers on topics including graph neural networks for
materials property prediction, physics-informed machine learning for
manufacturing processes, and Bayesian optimization for experimental
design. [If applicable: Previously served as [Role] at [Organization]
where they led development of ML models for thermal spray coating
optimization.] [PI Name] will devote [X]\% effort to this project.

\textbf{Co-Investigator:} [Co-PI Name] contributes [X] years of
expertise in thermal spray coating technology and surface engineering.
[Relevant credentials: PhD in materials science/mechanical engineering,
prior roles at coating companies or national laboratories, publications
on thermal spray process--structure--property relationships, experience
with APS, HVOF, and cold spray systems.] [Co-PI Name] will devote
[X]\% effort.

\textbf{Key Personnel:} A machine learning engineer (to be hired) will
support GNN model development and PINN implementation. A postdoctoral
researcher with expertise in equivariant neural networks and generative
models will contribute to digital twin development.

\textbf{Facilities and Equipment:} [Company Name] maintains
GPU-accelerated computing infrastructure (NVIDIA A100/H100 clusters)
for model training. A research collaboration with [Partner Institution]
provides access to APS and HVOF spray systems, DPV-2000 in-flight
particle diagnostics, SEM/EDS for coating characterization, and
mechanical testing (adhesion, hardness, wear).

\textbf{Relevant Prior Work:}
\begin{itemize}[nosep]
  \item [If applicable: Prior SBIR/STTR awards and their outcomes.]
  \item Demonstrated GNN property prediction achieving $R^2 > 0.85$ on
        materials benchmarks.
  \item Bayesian optimization campaigns achieving 10$\times$ efficiency
        gains over DOE for coating formulation.
  \item [Additional relevant experience in AI for manufacturing, digital
        twins, or coating technology.]
\end{itemize}

%% =======================================================================
\section{Cost and Schedule Summary}
\label{sec:cost}
%% =======================================================================

\begin{center}
\begin{tabular}{lrr}
\toprule
\textbf{Cost Category} & \textbf{Amount} & \textbf{\% of Total} \\
\midrule
Direct Labor (PI, Co-PI, ML Engineer) & \$132,000 & 55.0\% \\
Fringe Benefits (30\%) & \$39,600 & 16.5\% \\
Materials and Supplies & \$8,000 & 3.3\% \\
GPU Cloud Computing (AWS/Lambda) & \$15,000 & 6.3\% \\
Experimental Validation (spray trials) & \$18,000 & 7.5\% \\
Travel (NAVAIR/NAVSEA coordination) & \$6,000 & 2.5\% \\
Subcontractor ([Partner Institution]) & \$12,000 & 5.0\% \\
Indirect Costs & \$9,400 & 3.9\% \\
\midrule
\textbf{Total Phase I Cost} & \textbf{\$240,000} & \textbf{100.0\%} \\
\bottomrule
\end{tabular}
\end{center}

\textbf{Schedule:} 6-month period of performance. Key milestones at
Months 2 (dataset and pipeline), 3 (GNN Go/No-Go), 4 (digital twin
calibration), 5 (digital twin validation), and 6 (BO demonstration
and final report).

\textbf{Data Rights:} [Company Name] asserts limited rights in the
proprietary AI/ML algorithms, GNN architectures, and digital twin
models developed under this contract, consistent with DFARS
252.227-7018.

%% =======================================================================
\section*{References}
\label{sec:references}
%% =======================================================================

\begin{enumerate}[nosep,label={[\arabic*]}]
  \item Fauchais, P., Montavon, G., \& Bertrand, G. (2010). From
        powders to thermally sprayed coatings. \textit{Journal of
        Thermal Spray Technology}, 19(1), 56--80.

  \item Vaidya, A., Srinivasan, V., Streibl, T., Friis, M., Chi, W., \&
        Sampath, S. (2008). Process maps for plasma spraying of yttria-
        stabilized zirconia. \textit{Materials Science and Engineering: A},
        497(1--2), 239--253.

  \item Ganesh, I. (2013). A review on magnesium aluminate (MgAl$_2$O$_4$)
        spinel: synthesis, processing and applications. \textit{International
        Materials Reviews}, 58(2), 63--112.

  \item Raissi, M., Perdikaris, P., \& Karniadakis, G.~E. (2019).
        Physics-informed neural networks: A deep learning framework for
        solving forward and inverse problems involving nonlinear partial
        differential equations. \textit{Journal of Computational Physics},
        378, 686--707.

  \item Batchelor, T.~A.~A. et al. (2022). Bayesian optimization for
        accelerated materials discovery. \textit{Chemical Reviews},
        122(1), 156--232.

  \item Balandat, M. et al. (2020). BoTorch: A framework for efficient
        Monte-Carlo Bayesian optimization. \textit{Advances in Neural
        Information Processing Systems}, 33, 21524--21538.

  \item Skomski, D. et al. (2023). Graph neural networks for materials
        science and chemistry. \textit{Annual Review of Materials
        Research}, 53, 93--114.

  \item Bate, D., Rajak, P., \& Nakamura, A. (2022). Machine learning
        aided optimization of thermal spray processes: A review.
        \textit{Surface and Coatings Technology}, 446, 128880.

  \item Tran, R. et al. (2025). Open Molecules 2025 (OMol25): A large-scale
        dataset for molecular energy and force predictions.
        \textit{arXiv:2505.08762}.

  \item Chen, L., Yang, Y., \& Sampath, S. (2021). Data-driven process
        maps for thermal spray coating optimization.
        \textit{Journal of Thermal Spray Technology}, 30, 1988--2002.

  \item Assadi, H., Kreye, H., G\"artner, F., \& Klassen, T. (2016).
        Cold spraying---A materials perspective.
        \textit{Acta Materialia}, 116, 382--407.

  \item Elhoushi, S., Shaban, K., \& Ahmed, W.~H. (2024). Digital twin
        framework for thermal spray coating processes. \textit{Journal of
        Manufacturing Processes}, 98, 215--228.
\end{enumerate}

\end{document}
